\subsection{A Phenome-Wide Map of Gene-Trait Relevance}

% STRUCTURE: This section summarizes the full phenome-wide application as a resource.
% It answers: what does the field get from PIGEAN that it didn't have before?

Applying the Mouse+MSigDB model across {{NUM_TOTAL_TRAITS_MOUSE_MSIGDB:,}} traits ({{NUM_COMMON_TRAITS_MOUSE_MSIGDB:,}} common, {{NUM_RARE_TRAITS_MOUSE_MSIGDB:,}} rare) and {{NUM_TOTAL_GENE_ANNOTATIONS:,}} gene annotations, PIGEAN produces a dense probabilistic map of gene-trait relevance. At a posterior probability threshold of 90\%, we identify {{NUM_GENETRAIT_REL_PP90:,}} gene-trait relationships, an average of {{NUM_PER_TRAIT_PP90}} genes per trait and {{NUM_PER_GENE_PP90}} traits per gene. At a more inclusive threshold of 50\%, this expands to {{NUM_GENETRAIT_REL_PP50:,}} relationships ({{NUM_PER_TRAIT_PP50}} per trait, {{NUM_PER_GENE_PP50}} per gene).

% TODO: Add comparison to existing resources - how many gene-trait relationships
% does Open Targets Genetics have? GWAS Catalog? Orphanet? This would quantify
% the "orders of magnitude denser" claim from the abstract.

Each gene-trait relationship is accompanied by provenance identifying the specific SNP associations or gene records contributing to direct support and the specific genesets contributing to indirect support. Across all traits, PIGEAN evaluates {{NUM_GENESET_TRAIT_ASSOC_MILLIONS}} million geneset-trait associations, retaining only those that survive the spike-and-slab regularization as non-zero contributors to indirect support.

% TODO: Add a brief description of what users can do with this resource.
% - Query a gene: which traits is it relevant to, and through which annotations?
% - Query a trait: which genes are relevant, and which are novel (indirect-only)?
% - Query a geneset: for which traits does it carry non-zero effect?
% - Bring your own annotations: run PIGEAN with a new CRISPR screen or single-cell atlas

% FIGURE SUGGESTION: A summary figure showing (a) density of gene-trait map vs. existing
% resources, (b) distribution of direct vs indirect contributions across gene-trait pairs,
% (c) example queries for a well-known gene (e.g., PCSK9) and a well-known trait (e.g., T2D).

% ANALYSIS IDEA: Show that the resource captures gene-trait relationships not in any
% existing curated database. How many PP>90% relationships have no corresponding
% entry in Open Targets or GWAS Catalog?

% STRUCTURE: This section presents two anecdotes of how the resource can generate mechanistic
% hypotheses: (1) the HLA-DRB1/MS sex paradox case, and (2) cross-referencing PIGEAN gene set
% enrichments with an independent CRISPR screen in human adipocytes.

% STRUCTURE: First anecdote -- querying PIGEAN geneset enrichments for a known gene-trait
% pair to generate mechanistic hypotheses that connect siloed literatures.

As an illustration of how the resource can generate mechanistic hypotheses, we examined the PIGEAN geneset enrichments for \textit{HLA-DRB1} in multiple sclerosis (MS). A collaborator had observed that male MS patients carrying \textit{HLA-DRB1} risk variants exhibited worse inflammation biomarkers (serum GFAP and NfL) and worse clinical outcomes than female carriers of the same variants---a pattern consistent with the well-documented MS sex paradox in which women develop MS approximately three times more often but men progress faster once diagnosed \cite{voskuhl_sex_2020}. Querying the PIGEAN resource for the top geneset enrichments of \textit{HLA-DRB1} in MS across two model configurations (a smaller MSigDB-only model and a larger Mouse+MSigDB model) revealed convergent biology: cell adhesion molecules (KEGG CAMs, enrichment score 0.848 in the small model), lymphotoxin-alpha (LTA, 0.562), common variable immunodeficiency (0.571), CD40 signaling (0.150), TNF receptor superfamily member~7/CD27 (0.204), and multiple gene ontology terms related to leukocyte adhesion and T~cell activation.

% TODO: MS_HLA_DRB1_NUM_GENESETS_SMALL - Number of non-zero genesets in the small model run
% TODO: MS_HLA_DRB1_NUM_GENESETS_LARGE - Number of non-zero genesets in the large model run
% TODO: Add formal references for each key paper cited in this section to references.bib

These enrichments decompose the \textit{HLA-DRB1} effect into a three-stage mechanistic cascade supported by independent experimental evidence: (1)~enhanced MHC-II-mediated antigen presentation driving T~cell priming, consistent with reports that \textit{HLA-DRB1*15:01} carriers show constitutively elevated HLA-DR expression on cortical microglia and larger grey matter lesions \cite{enz_increased_2020, yates_association_2015}; (2)~TNF superfamily amplification via LTA, CD40, and CD27 signaling that promotes B~cell activation and meningeal tertiary lymphoid tissue (TLT) formation, as demonstrated experimentally by Bates et al.\ (2022) and Naouar et al.\ (2026) \cite{bates_lymphotoxin_2022, naouar_lymphotoxin_2026}; and (3)~cell adhesion molecule-mediated lymphocyte transmigration across the blood-brain barrier, the rate-limiting step for CNS immune infiltration and the therapeutic target of natalizumab. Critically, each stage of this cascade interacts with sex-specific biology: estrogen partially suppresses TNF/LTA production by activated lymphocytes, providing upstream braking in women that is absent in men; blood-brain barrier endothelial responses to inflammation are sexually dimorphic \cite{weber_cerebrovascular_2021}; and downstream neuroprotection through estrogen receptor signaling on astrocytes and oligodendrocytes buffers women against the cortical damage driven by meningeal TLTs \cite{spence_neuroprotection_2011}. Men carrying \textit{HLA-DRB1} risk variants thus face a ``double hit'': the immune cascade runs with less hormonal restraint while the CNS neuroprotective buffer is weaker, consistent with the recent finding that \textit{HLA-DRB1*15:01} is associated with reduced longevity specifically in males \cite{mendoza_mejia_hla_2025}. While individual components of this model have been characterized in separate literatures---sex differences in MS progression, HLA-DRB1 effects on cortical pathology, and LTA-driven TLT formation---the integrated chain linking genotype to sex-differential severity through these specific intermediate mechanisms had not been previously proposed. The PIGEAN geneset enrichments independently recovered these nodes from GWAS data alone, illustrating how the resource can bridge siloed experimental literatures to generate testable mechanistic hypotheses.

% MISSING EVIDENCE: The sex-stratified analysis of HLA-DRB1 effects on GFAP/NfL
% would be strengthened by including the specific effect sizes from the collaborator's
% analysis once available. Consider adding variables:
% TODO: MS_HLA_DRB1_MALE_GFAP_BETA - Effect of DRB1 on GFAP in males
% TODO: MS_HLA_DRB1_FEMALE_GFAP_BETA - Effect of DRB1 on GFAP in females
% TODO: MS_HLA_DRB1_SEX_INTERACTION_P - P-value for sex interaction term

% ANALYSIS IDEA: Run PIGEAN on MS GWAS stratified by sex (if sex-stratified
% summary statistics are available) to test whether the geneset enrichment profile
% itself differs between male and female MS. This would directly test whether
% the model captures sex-differential biology at the geneset level.

% STRUCTURE: Second anecdote -- cross-referencing PIGEAN indirect support with an
% independent CRISPR screen to identify genes whose known function predicts a novel
% relationship with adipocyte biology.

\subsubsection{Cross-referencing indirect support with functional screens}

We next cross-referenced PIGEAN predictions with an independent high-content CRISPR screen in human adipocytes \cite{flannick_adipocyte_2024}. The screen profiled {{CRISPR_N_GENES}} genes across imaging phenotypes including lipid droplet morphology (BODIPY), actin cytoskeleton texture, and nuclear morphology (DAPI), clustering genes into nine factors. We queried the PIGEAN resource for these {{CRISPR_N_GENES}} genes across {{CRISPR_N_PHENOTYPES_MSIGDB}} curated metabolic, glycemic, hepatic, and anthropometric traits and classified each gene by its evidence type: \textit{understood mechanism} (high indirect and high direct support), \textit{novel GWAS discovery} (high direct, low indirect), or \textit{indirect-only} (high indirect support, low direct support). Of the screened genes, {{CRISPR_N_HITS}} ({{CRISPR_PCT_HITS}}\%) had at least one PIGEAN hit (combined support $\geq$ 2.0), including {{CRISPR_N_KNOWN}} of 42 known adipose-metabolism genes (\textit{PPARG} on rare diabetes, indirect support = 6.87; \textit{LPL} on TG:HDL ratio, indirect support = 6.61).

% TODO: CRISPR_N_GENES - Number of unique CRISPR screen genes (230)
% TODO: CRISPR_N_PHENOTYPES_MSIGDB - Number of metabolic phenotypes in msigdb model (537)
% TODO: CRISPR_N_HITS - Number of genes with at least one hit (114)
% TODO: CRISPR_PCT_HITS - Percentage of genes with hits (49.6)
% TODO: CRISPR_N_KNOWN - Number of known adipose genes with hits (40)
% TODO: CRISPR_N_NOVEL_INDIRECT - Number of novel indirect-only genes (5)

{{CRISPR_N_NOVEL_INDIRECT}} genes had high indirect support but low direct support for metabolic traits while also appearing in the CRISPR screen---genes whose pathway biology predicts metabolic relevance independently of any GWAS signal. \textit{EPS15} received the highest indirect support (3.64 for HbA1c adjusted for BMI; direct support = 1.30), driven by gene sets for decreased circulating insulin ($\beta$ = 1.17) and increased circulating glucose ($\beta$ = 1.02). EPS15 is an endocytic adaptor protein required for clathrin-coated pit assembly \cite{vanbergen_eps15_1995}, and the CRISPR screen assigned it to the lipid droplet morphology factor. In adipocytes, insulin-stimulated GLUT4 translocation to the plasma membrane depends on clathrin-dependent endocytic recycling \cite{leto_regulation_2012}; disruption of EPS15-dependent endocytosis would impair both insulin receptor internalization and GLUT4 re-exocytosis, consistent with the insulin/glucose gene set signal.

\textit{ADRA1B}, encoding the alpha-1B adrenergic receptor, received indirect support of 3.22 for fasting insulin adjusted for BMI (direct support = 1.11), driven by gene sets for abnormal glucose homeostasis ($\beta$ = 1.16), increased circulating insulin ($\beta$ = 1.10), and impaired glucose tolerance ($\beta$ = 1.07). The CRISPR screen placed it in the same lipid droplet morphology factor as \textit{EPS15}. Alpha-2 adrenergic receptors are established regulators of catecholamine-stimulated lipolysis in white adipose tissue \cite{lafontan_fat_2009, bartness_sympathetic_2014}, but a role for the alpha-1B subtype in adipocyte insulin sensitivity has not been reported. The convergence of a lipid droplet morphology phenotype in the screen with insulin/glucose gene set enrichments in PIGEAN suggests ADRA1B may modulate adipocyte insulin signaling through a mechanism distinct from classical lipolytic control.

\textit{RNF213}, an AAA+ ATPase and E3 ubiquitin ligase known as the major susceptibility gene for Moyamoya disease \cite{sugihara_rnf213_2019}, received indirect support of 2.30 for 2-hour glucose adjusted for BMI (direct support = 1.12), driven by gene sets for increased circulating glucose ($\beta$ = 1.32) and impaired glucose tolerance ($\beta$ = 1.12). The CRISPR screen assigned it to the actin/metabolic traits factor. Recent work has shown that RNF213 ubiquitinates proteins on the surface of intracellular lipid droplets \cite{ahel_rnf213_2023}, and the combination of this lipid droplet remodeling activity with PIGEAN-predicted glucose homeostasis relevance suggests RNF213 may regulate lipid droplet turnover in adipocytes through ubiquitin-dependent quality control. None of these three genes would have been prioritized by direct genetic association alone (all direct support $<$ 1.3).

% FIGURE SUGGESTION: A 2D scatter plot of indirect vs direct support for all 74 novel CRISPR
% genes, colored by evidence bucket (indirect-only, understood mechanism, novel GWAS,
% modest). Annotate the indirect-only genes. This would visually communicate the
% three-bucket framework and highlight the indirect-only quadrant.
